\documentclass{rvcepaper}
\begin{document}

% ---------- HS351TA Entrepreneurship & IPR, CIE - I, 13th April 2026
\begin{iempaper}
\paperinfo{shownote=0,date={13\textsuperscript{th} April 2026},exam={CIE -- I},maxmarks={10 + 50},sem={VI},prog={UG},duration={30 + 90 Min},
 title={Entrepreneurship \& Intellectual Property Rights},code={HS351TA}}
\iemhead
{\small Note:}\par\vspace{-1pt}
\hspace*{0.5cm}{\footnotesize 1. \hspace{0.1cm}Answer all the Questions.}\par\vspace{3pt}
\begin{cietable}{M}{BT}{CO}
\qpart{Part -- A}
\q{1.}{Define Certification Trade Mark.}{02}{2}{5}
\q{2.}{The Indian government allows a local company to produce a patented medicine during a health emergency without the owner's consent. What is this called under patent law?}{02}{2}{5}
\q{3.}{A company tries to register the word ``Laptop'' as its trademark for selling computers. Can this mark be registered?}{02}{2}{5}
\q{4.}{Dr. Meena discovers a new mathematical formula for predicting weather patterns. Can this discovery be patented?}{02}{2}{5}
\q{5.}{Analyze how the networking skills influence entrepreneurial success.}{02}{2}{1}
\qpart{Part -- B}
\q{1}{Critically assess the statement ``Entrepreneurs never fail'' and explain how failure is an integral part of entrepreneurial growth.}{10}{2}{1}
\q{2}{Develop a case study of a successful entrepreneurial venture that utilized a unique method for identifying its opportunity, explaining the approach and its results.}{10}{2}{1}
\q{3}{An international company develops a method to extract active compounds from a plant traditionally used in Ayurveda, claiming higher efficiency and stability. Can this method be patented, and how does it relate to the protection of traditional knowledge?}{10}{3}{5}
\q{4}{Discuss the concept of patent infringement in detail. Include various types of infringement with relevant case studies to illustrate the practical application of patent law.}{10}{3}{4}
\q{5}{Explain the concepts of Patent Search and Patent Drafting. Discuss their importance in the process of patent protection.}{10}{3}{5}
\end{cietable}
\end{iempaper}

% ---------- HS361TA Entrepreneurship & IPR, CIE - I, 28th April 2025
\begin{iempaper}
\paperinfo{shownote=0,date={28\textsuperscript{th} April 2025},exam={CIE -- I},maxmarks={10 + 50},sem={VI},prog={UG},duration={30 + 90 Min},
 title={Entrepreneurship \& Intellectual Property Rights},code={HS361TA}}
\iemhead
{\small Note:}\par\vspace{-1pt}
\hspace*{0.5cm}{\footnotesize 1. \hspace{0.1cm}Answer all the Questions.}\par\vspace{3pt}
\begin{cietable}{M}{BT}{CO}
\qpart{Part -- A}
\q{1.}{What is the primary difference between an Innovative Entrepreneur and an imitative entrepreneur?}{1}{2}{1}
\q{2.}{Why do Fabian entrepreneurs tend to be less successful in fast-paced industries?}{1}{2}{1}
\q{3.}{What is the main reason for the failure of many new entrepreneurial ventures?}{1}{2}{1}
\q{4.}{Why is risk-taking considered an essential characteristic of an entrepreneur?}{1}{2}{1}
\q{5.}{What emerging trend in entrepreneurship has been accelerated by digital technologies?}{1}{2}{1}
\q{6.}{What is the primary difference between ``entrepreneurship'' and ``intrapreneurship''?}{1}{2}{1}
\q{7.}{What is the primary purpose of intellectual property (IP) laws?}{1}{2}{4}
\q{8.}{If an inventor publishes their invention before filing a patent, what impact does it have on the patentability?}{1}{2}{4}
\q{9.}{Which of the following can be patented: A new business model, a new type of software algorithm, or a new pharmaceutical drug?}{1}{2}{4}
\q{10.}{Can a product that is patented in one country be sold without a patent in another country?}{1}{2}{5}
\qpart{Part -- B}
\q{1a.}{Define Entrepreneurship and explain its scope in the context of engineering. Discuss the various areas where entrepreneurship plays a key role in engineering industries.}{5}{2}{1}
\q{b.}{What are the characteristics and traits of successful entrepreneurs? Discuss how these qualities are essential for engineering entrepreneurs to thrive in a competitive and rapidly changing market.}{5}{2}{1}
\q{2.}{\textbf{FitTrack Innovations} was a wearable tech startup launched in 2017, aiming to develop a fitness tracker focused on personalized health insights. Despite having a unique product, the startup struggled with low sales and negative user feedback. The key issues were poor marketing, unclear product positioning, and a lack of customer engagement.\newline
After the initial failure, the founders decided to pivot. They conducted detailed market research, redesigned the product to focus on a more specific niche---tracking chronic health conditions---and revamped their marketing strategy. With a clearer value proposition and targeted audience, \textbf{FitTrack Innovations} relaunched and gained significant traction in the health tech market. The company later secured venture capital and expanded its product range.
\begin{abc}\item What do you think was the most critical factor in FitTrack Innovations' failure, and how could the founders have addressed it earlier to avoid the setback?\item What is a pivot in the context of startups, and why did FitTrack Innovations decide to pivot their business model?\item How would you evaluate the effectiveness of the founders' decision to pivot? What risks might they have faced, and what alternatives could they have considered?\item What are the key factors that determine the success or failure of a product launch for a startup?\item How does securing venture capital help a struggling startup, and how did it impact the recovery of FitTrack Innovations?\end{abc}}{10}{3}{2}
\q{3.}{`SunRay' Innovations is an engineering startup that has developed a cutting-edge solar powered water purification system aimed at providing clean drinking water to rural areas. The technology is innovative and has the potential to improve the quality of life for many communities. However, `SunRay' Innovations is struggling to secure funding and has limited market reach due to strong competition from larger companies with more resources.
\begin{abc}\item Discuss how `SunRay' Innovations demonstrates the key principles of entrepreneurship through its innovation.\item Identify potential entrepreneurial opportunities for `SunRay' Innovations to overcome its challenges and scale its business.\item How does `SunRay' Innovations' focus on sustainable technology reflect the broader role of entrepreneurship in addressing social and environmental challenges?\item How can the company effectively manage risks associated with expansion?\end{abc}}{10}{3}{2}
\q{4a.}{Define Intellectual Property (IP) and explain its importance in the modern economy. Discuss the different types of intellectual property and their significance for businesses.}{5}{2}{4}
\q{b.}{Provide an overview of the patent application procedure. Discuss the key steps involved in securing a patent and the timelines for the process.}{5}{2}{4}
\q{5.}{InnovTech, a startup based in India, has developed an innovative device that converts waste plastic into reusable energy. The company has already filed for a patent for their invention. However, after conducting a thorough patent search, they found that a similar invention had already been patented in the United States. InnovTech believes their technology is more efficient and has a better potential for scalability.\newline
Now, InnovTech is considering commercializing their invention both locally and globally. They are also exploring the possibility of licensing their patent to larger corporations or transferring the patent rights.
\begin{abc}\item Discuss the importance of conducting a patent search in the early stages of product development. How does this process help InnovTech in ensuring the uniqueness of their invention?\item Discuss the steps InnovTech should take to protect its invention and ensure its patentability.\item If InnovTech decides to pursue international protection for their invention, what steps must they follow to ensure global patent protection?\item What options does InnovTech have for transferring or commercializing their patent? What legal remedies could they seek if their patent rights are infringed upon?\end{abc}}{10}{3}{5}
\end{cietable}
\end{iempaper}

% ---------- HS351TA Entrepreneurship & IPR, CIE - I, 25th November 2024
\begin{iempaperB}
\paperinfo{ayear={2024-2025 (ODD Sem)},date={25\textsuperscript{th} November 2024},code={HS351TA},sem={V Semester},exam={CIE -- I},maxmarks={10 + 50},duration={120 Mins},
 title={ENTREPRENEURSHIP AND INTELLECTUAL PROPERTY RIGHTS}}
\iemheadB
\begin{cietable}{M}{BT}{CO}
\qpart{PART -- A}
\q{1}{Define Compulsory license under the patent law.}{02}{1}{4}
\q{2}{List any two characteristics of successful entrepreneurs.}{02}{2}{1}
\q{3}{Difference between Intrapreneur \& Entrepreneur.}{02}{1}{1}
\q{4}{List any two non-patentable inventions under the Indian patent Act.}{02}{2}{2}
\q{5}{List any two myths of entrepreneurs.}{02}{2}{1}
\qpart{PART -- B}
\q{1a}{A software company develops an AI tool for legal research and applies for a patent in India. Discuss the challenges it might face under Indian patent law.}{05}{3}{3}
\q{b}{Explain the difference between a product patent and a process patent with an examples.}{05}{1}{4}
\q{2}{Create a commercialization strategy for a patent related to renewable energy technology.}{10}{3}{3}
\q{3}{Propose a framework for enhancing Traditional Knowledge protection under international Intellectual Property laws.}{10}{2}{4}
\q{4}{Create a case study on how perseverance as a characteristic helped a famous entrepreneur achieve success.}{10}{3}{1}
\q{5}{Explain how entrepreneurship fosters innovation in an economy with one case example.}{10}{2}{1}
\end{cietable}
\btfoot{|l|l|l|*{11}{C{0.8cm}|}}{\hline
\multirow{2}{*}{\shortstack{Marks\\Distribution}} & \multicolumn{2}{c|}{Particulars} & CO1 & CO2 & CO3 & CO4 & CO5 & L1 & L2 & L3 & L4 & L5 & L6\\\cline{2-14}
 & Max & Quiz & 04 & 02 & -- & 04 & -- & 02 & 02 & 06 & -- & -- & --\\\cline{3-14}
 & Marks & Test & 20 & 10 & 15 & 15 & -- & 05 & 10 & 25 & -- & -- & --\\\hline}
\end{iempaperB}

\end{document}
